\section{Future work}
\label{future-work}
\subsection{Features}
\subsubsection{State management}

The host environment presented in this paper makes ILD mostly (e.g.
\paraphrase{not counting} \ild{gensym}) purely-functional and \paraphrase{makes
no effort} to allow the user to manage \paraphrase{state}. However, for
general-purpose applications, ILD must be embedded in an impure environment
that allows state management and IO.

Other functional languages that feature fully-immutable data employ different
methods of handling state, IO and other side effects:
\begin{itemize}
    \item Clojure fetures immutable core data structures but allows mutable
    state management via object handles such as \code{atom}/\code{ref}:
    the handle itself is immutable, but operations associated with the handle
    may mutate/replace its underlying value~\citneeded. IO operations in clojure
    also work via similar references. Haskell has a similar
    facillity (e.g. \code{IORef}, \code{MVar})~\citneeded
    \item Erlang, LFE and other languages on the BEAM virtual machine employ
    processes and message passing as primitives and manage state via an
    actor pattern~\citneeded
    \item \paraphrase{FRP?}
\end{itemize}
Embedding ILD in environments with similar practical features is worth exploring.

\subsubsection{Quasiquote}
ILD's \ildsf{quote} special form yields the given program without
evaluating it. One use for \ildsf{quote} is for a macro function to
return a literal \paraphrase{chunk} of code as its result. However,
plain \ildsf{quote} does not allow dynamic values to be easily inserted
inside the returned expression. To remedy this, many LISP-like languages
employ a \emph{quasiquote} form that looks for special \emph{unquote} markers
inside the given expression and evaluates the unquoted subexpressions before
splicing them into the result.~\citneeded A similar mechanism can be implemented
in ILD using the metaprogramming techniques shown in this paper.
\subsection{Optimisations}
\subsubsection{Partial evaluation}
In a naïve implementation of ILD's interpreter, macros are expanded runtime, as encountered.
This strategy incurs a significant cost, possibly exponential in the depth
of encountered nested macro invocations, since the body of each macro call
is re-evaluated each time the surrounding expression is evaluated. The naïve
ILD interpreter takes a whole second to evaluate the program from \cref{example-module-fact}
on a modern machine due to exponential macro call reevaluation.
Nevertheless, performance can be greatly improved by evaluating macros ahead of time.
We argue that placing the boundary between code evaluated ahead of time and code
evaluated at runtime should not be at the distinction between ``macro expansion'' and
``code execution'', like LISP does, but rather between ``code dependant on data known ahead of time''
(which includes most, but not all, macros) and ``code dependant on runtime data''.
We aim to explore this topic in further research by applying partial evaluation
as an optimisation pass in ILD, similar to how this has been done in FExpr-based languages
like Kraken \cite{fexpr-pe}. In addition to eliminating the overhead
of runtime macro semantics, PE would also optimise other, non-macro code that does not
depend on user input.
\subsubsection{Tail call optimisation}
Languages such as Scheme, Lisp or Racket support expressing efficient iterative
algorithms as tail-recursive functions, which get optimised into \paraphrase{trampolined}
tail calls by the interpreter.~\citneeded A TCO-enabled ILD interpreter is worth
exploring \paraphrase{for similar reasons.}
\subsubsection{Module caching}
The simple module loader program as presented in this paper evaluates a module every
time it appears in the dependency graph of the root module being loaded.
Due to lack of mutable module state, this is only an efficiency problem. A less
naïve module loader could traverse the dependency graph ahead of time and
resolve module imports without duplicates

